\documentclass[10pt,a4paper]{amsart}
\usepackage[margin=19mm]{geometry}
\usepackage{amsmath,amssymb,mathtools}
\usepackage[scr=rsfs,cal=euler]{mathalfa}
\usepackage{xfrac}
\usepackage[unicode,hidelinks,bookmarks=false]{hyperref}
\newcommand{\ie}{i.e.\ }
\newcommand{\cf}{cf.\ }
\newcommand{\resp}{respectively\ }
\newcommand{\Subsection}{\subsection}
\providecommand{\nobreakdash}{\nobreak\hbox{-}}
\setlength{\emergencystretch}{2em}
\multlinegap0pt
\usepackage{bm}
\usepackage{bbm}
\usepackage{dsfont}
\usepackage{xspace}
\usepackage{cancel}
\usepackage{mathtools}
\usepackage{amsthm}
\makeatletter
\@ifundefined{theorem}{\newtheorem{theorem}{Theorem}}{}
\@ifundefined{claim}{\newtheorem{claim}[theorem]{Claim}}{}
\@ifundefined{corollary}{\newtheorem{corollary}[theorem]{Corollary}}{}
\@ifundefined{proposition}{\newtheorem{proposition}[theorem]{Proposition}}{}
\makeatother
\title[A combinatorial approach to nonlinear spectral gaps]{A combinatorial approach to nonlinear spectral gaps}
\author{Dylan J. Altschuler, Pandelis Dodos, Konstantin Tikhomirov, Konstantinos Tyros}
\date{Source: arXiv:2410.04394v3 (CC BY 4.0)}
\begin{document}
\maketitle

\noindent\emph{Source and licence.} A combinatorial approach to nonlinear spectral gaps. Version arXiv:2410.04394v3 (CC BY 4.0), distributed under the Creative Commons Attribution 4.0 International licence, as recorded in the arXiv licence metadata for this exact version and shown on the abstract page https://arxiv.org/abs/2410.04394v3. Full rights notice: licenses/arxiv-2410.04394-CC-BY-4.0.txt.\par\smallskip

\noindent\textit{Editorial scope.} Counted unit: the Proposition whose source label is \texttt{p3-pap1} in the article (source0.tex lines 1539--1556, source label \texttt{p3-pap1}), delivered together with the article's own complete original proof of it (source0.tex lines 1557--1706, one \texttt{proof} environment of 150 lines). No mathematics is written, repaired or re-derived: statement and proof are the article's own text line for line. Required context retained uncounted: c2-pap1 (lines 1477--1487); eq:new\_02 (lines 1503--1506); eq:new\_04 (lines 1508--1510); cor\_new\_base\_step (lines 1524--1529). Every definition, display and label of those lines is retained unchanged. Omitted from this package and not redistributed: 1--1476, 1488--1502, 1507--1507, 1511--1523, 1530--1538, 1707--2779. Nothing inside a retained excerpt is omitted or altered: every retained body is delivered exactly as the article's lines are written, so no omission receipt of any kind accompanies this package. Citations: the retained excerpts carry no citation command at all, so no bibliography entry of the article is transcribed and no reference is redistributed. Reference adaptation: none was needed and none was made; every reference inside the delivered unit resolves inside it, so no unresolved reference or citation is produced. Printed slips kept literal: no printed word, symbol, hypothesis or number of the counted statement or of its proof was altered. Layout and class: the article's own class is \texttt{amsart} with packages including graphicx, xcolor, fullpage, amsmath, amsthm, amssymb, latexsym, color, hyperref, lipsum, enumitem, bbm, tikz, subfigure; the wrapper is the lane's pinned \texttt{amsart} editorial wrapper at 10pt on A4, which restates the theorem-like environments the retained text needs and adds the article's own macro definitions that the retained text needs (none), with extra wrapper packages (bm, bbm, dsfont, xspace, cancel, mathtools). The delivered numbering is the unit's own. Assets: no retained span contains a figure environment, a graphics-inclusion command, a PDF-page inclusion, a table or a TikZ picture; no image file is copied into this package, the package declares no asset dependency and no image file is redistributed with it. Licence: version arXiv:2410.04394v3 of this article is distributed under the Creative Commons Attribution 4.0 International licence, as recorded in the arXiv licence metadata for this exact version and shown on the abstract page https://arxiv.org/abs/2410.04394v3. A full rights notice is delivered at licenses/arxiv-2410.04394-CC-BY-4.0.txt. Characters: every retained excerpt is pure ASCII, so the delivered engine, XeLaTeX with its default Latin Modern text font, reports no missing character.
\begin{corollary} \label{c2-pap1}
Let $S\subseteq [n]$ be nonempty, and let $\ell$ be a positive integer such that
$|S|d(d-1)^{\ell-2}\leqslant\frac{n}{6}$ if $\ell\geqslant 2$, and $|S|\leqslant \frac{n}{6}$ if $\ell=1$.
Also let $0<\theta<1$, and let $\pi\in\mathcal{M}^S_{\ell-1}$. Then,
\begin{align*}
\mathbb{P}\Big[ \big|\Delta_{\mathbf{G}}(S,\ell)\big|\geqslant & \
\theta (d-1) \big|\Delta_{G(\pi)}(S,\ell-1)\big| \, \Big| \, \mathcal{M}_{\pi} \Big] \geqslant \\
& 1 - \Big(\frac{4e}{1-\theta}\cdot \frac{d(d-1)^{\ell-1}|S|}{n}
\Big)^{\frac{1-\theta}{2} (d-1) |\Delta_{G(\pi)}(S,\ell-1)|}.
\end{align*}
\end{corollary}

\begin{equation}\label{eq:new_02}
L_0 := \big\lfloor(7/15)10^8\big\rfloor+1 \ \ \text{ and } \ \
K := \frac{d-2.1\sqrt{d-1}}{2}\Big( 1.5-1.05\,\frac{\sqrt{d-1}}{d} \Big)^{L_0-1}
\end{equation}

\begin{equation}\label{eq:new_04}
K\geqslant \frac{28}{0.008}=\frac{7}{2}10^3.
\end{equation}

\begin{corollary}\label{cor_new_base_step}
Let $S\subseteq [n]$ be nonempty with $d(d-1)^{L_0-2}|S|\leqslant\frac{n}{4}$, where $L_0$ is as in \eqref{eq:new_02}.
Also let $\pi_0\in\mathcal{M}^S_{L_0-1}$ such that $\mathcal{M}_{\pi_0}\cap\mathcal{E}\neq\emptyset$. Then,
\[\mathbb{P}\Big[ \big|\Delta_{\mathbf{G}}(S,L_0)\big|\geqslant \frac{K}{2} |S| \, \Big| \, \mathcal{M}_{\pi_0} \Big] \geqslant
1 - \Big(8e\, \frac{d(d-1)^{L_0-1}|S|}{n} \Big)^{\frac{K}{4} |S|}.\]
\end{corollary}

\begin{proposition} \label{p3-pap1}
Set
\begin{equation} \label{e-pap1-p1}
\eta=\eta(d):= \frac{1}{12^{2} e^{3} (d-1)^{2L_0+2}},
\end{equation}
where $L_0$ is as in \eqref{eq:new_02}.
Let $S\subseteq [n]$ be nonempty with $|S|\leqslant \eta n$,
let $\pi_0\in\mathcal{M}^S_{L_0-1}$ such that $\mathcal{M}_{\pi_0}\cap\mathcal{E}\neq\emptyset$, and set
\begin{equation} \label{e-pap-p2}
T:= \Big\lfloor \log_{d-1}\Big(\frac{\eta n}{|S|}\Big)\Big\rfloor.
\end{equation}
Then,
\begin{equation} \label{eqkt22}
\mathbb{P} \Big[ \big|\partial B_{\mathbf{G}}(S,L_0+\ell)\big| \geqslant
28  (d-1)^\ell |S| \text{ for all } \ell\in [T]  \, \Big| \, \mathcal{M}_{\pi_0} \Big]
 \geqslant 1- \Big(\frac{|S|}{en}\Big)^{3|S|}.
\end{equation}
\end{proposition}

\begin{proof}
Set $r:=\frac{1}{\sqrt{2}}$ and
\begin{equation} \label{e-pap-p3}
\begin{cases}
\delta_\ell := r^{\ell} & \text{if } \, \ell\in \big\{1,\dots,\lfloor 2T/3 \rfloor\big\}, \\
\delta_\ell :=  r^{1+2T-2\ell} & \text{if } \, \ell\in \big\{ \lfloor 2T/3 \rfloor+1,\dots,T\big\},
\end{cases}
\ \ \ \
\begin{cases}
\alpha_0:=\frac{1}{2}, \\
\alpha_\ell := \prod_{t=0}^{\ell}(1-\delta_t) & \text{if } \, \ell\in \{1,\dots,T\}.
\end{cases}
\end{equation}
Observe that $(\alpha_\ell)_{\ell=0}^T$ is decreasing, and
\begin{equation} \label{"douleia"}
\alpha_T\geqslant \frac{1}{2}\Big(\prod_{t=1}^{\infty}(1-r^t)\Big)
\Big(\prod_{t=0}^{\infty}(1-r r^{2t})\Big) \geqslant 0.008.
\end{equation}
Since $\Delta_{\mathbf{G}}(S,\ell)$ is a subset of $\partial B_{\mathbf{G}}(S,\ell)$ everywhere
in our probability space, by \eqref{eq:new_04}, the desired estimate \eqref{eqkt22} will follow once we show that
\begin{equation} \label{eqkt23}
\mathbb{P}\Big[ \big|\Delta_{\mathbf{G}}(S,L_0+\ell)\big| \geqslant
0.008 K (d-1)^\ell |S| \text{ for all } \ell\in [T]\Big]\geqslant 1- \Big(\frac{|S|}{en}\Big)^{3|S|},
\end{equation}
where $K$ is as in \eqref{eq:new_02}. For every nonnegative integer $\ell$, set
\begin{align*}
\widetilde{\mathcal{M}}^S_\ell & := \big\{ \pi\in\mathcal{M}^S_{L_0+\ell}\colon
\mathcal{M}_{\pi_0}\subseteq\mathcal{M}_\pi\text{ and }
\big| \Delta_{G(\pi)}(S,L_0+\ell)\big| \geqslant \alpha_\ell K (d-1)^\ell|S|\big\}, \\
\widehat{\mathcal{M}}^S_\ell & := \big\{\pi\in\mathcal{M}^S_{L_0+\ell} \colon
\mathcal{M}_{\pi_0}\subseteq\mathcal{M}_\pi\text{ and }
\big|\Delta_{G(\pi)}(S,L_0+t)\big| \geqslant \alpha_t K (d-1)^t|S| \text{ for all } t\in\{0,\dots,\ell\}\big\}.
\end{align*}
%{\color{red}Condition $\mathcal{M}_{\pi_0}\subseteq\mathcal{M}_\pi$ can be substituted equivalently by
%$\pi_0\subseteq\pi$ or $\mathcal{M}_{\pi_0}\cap\mathcal{M}_\pi\neq\emptyset$. Although $\pi_0\subseteq\pi$ is
%kind of more basic, the one I put there describes more directly what is going on with the events of the probability space. Just mentioning %it Prof. Pandelis to pick whatever you prefer.}
Clearly, $\widehat{\mathcal{M}}^S_\ell\subseteq\widetilde{\mathcal{M}}^S_\ell$
for all $\ell\in \{0,\ldots,T\}$. Moreover, by Corollary \ref{cor_new_base_step}, the choice of $\alpha_0$ and the fact that $\frac{d}{d-1}\leqslant \frac{3}{2}$, we have
\[\mathbb{P}\Big[ \big|\Delta_{\mathbf{G}}(S,L_0)\big|\geqslant \alpha_0 K |S| \, \Big| \, \mathcal{M}_{\pi_0} \Big] \geqslant
  1 - \Big(12e\, \frac{(d-1)^{L_0}|S|}{n} \Big)^{\frac{K}{4} |S|}.\]
On the other hand, by Corollary \ref{c2-pap1} and the fact that $\frac{d}{d-1}\leqslant \frac{3}{2}$,
for every positive integer $\ell$ and every $\pi\in \widetilde{\mathcal{M}}^S_{\ell-1}$, we have
\[ \mathbb{P}\Big[ \big|\Delta_{\mathbf{G}}(S,L_0+\ell)\big| \geqslant  \alpha_\ell K (d-1)^\ell |S| \, \Big| \,
\mathcal{M}_{\pi} \Big]
\geqslant 1 - \Big(\frac{6e}{\delta_\ell} \cdot
\frac{(d-1)^{L_0+\ell}|S|}{n}\Big)^{\frac{\delta_\ell}{2} \alpha_{\ell-1}K (d-1)^{\ell} |S|}.\]
Next, observe that for every $\ell\in\{0,\dots,T\}$, every distinct $\pi_1,\pi_2\in \mathcal{M}^S_{L_0+\ell}$ are incomparable
under inclusion which, in turn, implies that $\mathcal{M}_{\pi_1}\cap\mathcal{M}_{\pi_2}=\emptyset$;
consequently, we obtain that
\[\mathbb{P}\bigg[ \bigcup_{\pi\in \widehat{\mathcal{M}}^S_0} \mathcal{M}_{\pi} \, \bigg| \, \mathcal{M}_{\pi_0} \bigg] \geqslant
  1 - \Big(12e\, \frac{(d-1)^{L_0}|S|}{n}
\Big)^{\frac{K}{4} |S|}\]
and, for every $\ell\in[T]$,
\[\mathbb{P}\bigg[ \bigcup_{\pi\in \widehat{\mathcal{M}}^S_\ell} \mathcal{M}_{\pi} \, \bigg| \,
\bigcup_{\pi\in \widehat{\mathcal{M}}^S_{\ell-1}} \mathcal{M}_{\pi} \bigg]
\geqslant 1 - \Big(\frac{6e}{\delta_\ell} \cdot
\frac{(d-1)^{L_0+\ell}|S|}{n}\Big)^{\frac{\delta_\ell}{2} \alpha_{\ell-1}K (d-1)^{\ell} |S|}.\]
Hence, setting $\widetilde{\mathcal{M}}^S_{-1} = \widehat{\mathcal{M}}^S_{-1} := \{\pi_0\}$,
by a union bound  we have
\begin{align}
\mathbb{P}\bigg[  \bigcap_{\ell=0}^T  \Big[ \big|\Delta_{\mathbf{G}}(S, & L_0+\ell)\big| \geqslant
0.008 K (d-1)^\ell |S|  \big]  \, \bigg| \, \mathcal{M}_{\pi_0}  \bigg] \label{e-pap-p4} \\
&\geqslant
\mathbb{P}\bigg[ \bigcap_{\ell=0}^T \Big[ \big|\Delta_{\mathbf{G}}(S,L_0+\ell)\big| \geqslant
\alpha_\ell K(d-1)^\ell |S| \Big]  \, \bigg| \, \mathcal{M}_{\pi_0} \bigg] \nonumber\\
& = \mathbb{P}\bigg[ \bigcap_{\ell=1}^T \bigcup_{\pi\in \widehat{\mathcal{M}}^S_\ell} \mathcal{M}_{\pi} \, \bigg| \, \mathcal{M}_{\pi_0} \bigg] =
\prod_{\ell=1}^{T} \mathbb{P}\bigg[ \bigcup_{\pi\in\widehat{\mathcal{M}}^S_\ell} \mathcal{M}_{\pi} \, \bigg| \,
\bigcup_{\pi\in \widehat{\mathcal{M}}^S_{\ell-1}} \mathcal{M}_{\pi} \bigg] \nonumber \\
& \geqslant 1 -\Big(12e\, \frac{(d-1)^{L_0}|S|}{n} \Big)^{\frac{K}{4} |S|}
- \sum_{\ell=1}^{T}  \Big(\frac{6e}{\delta_\ell} \cdot
\frac{(d-1)^{L_0+\ell}|S|}{n}\Big)^{\frac{\delta_\ell}{2} \alpha_{\ell-1}K (d-1)^{\ell} |S|}. \nonumber
\end{align}
Thus, in order to prove \eqref{eqkt23}, it suffices to show that the quantities subtracted  in the right-hand-side of \eqref{e-pap-p4}
add up to at most $\binom{|S|}{e n}^{3|S|}$. This, in turn, would follow once we show that
\begin{equation}\label{eq:new_03}
\bigg(12 e^2 (d-1)^{L_0} \Big(\frac{|S|}{en}\Big)^{1-\frac{12}{K}} \bigg)^{\frac{K}{4}|S|} \leqslant \frac{1}{2}
\end{equation}
and, for every $\ell\in[T]$,
\begin{equation} \label{e-pap-p5}
\bigg(6e^2(d-1)^{L_0}\frac{(d-1)^{\ell}}{\delta_\ell} \cdot
\Big(\frac{|S|}{en}\Big)^{1-\frac{6}{\delta_\ell\alpha_{\ell-1}K(d-1)^{\ell}}}
\bigg)^{\frac{\delta_\ell}{2} \alpha_{\ell-1}K (d-1)^{\ell} |S|}
\leqslant \frac{1}{2^{\ell+1}}.
\end{equation}

First we argue for \eqref{eq:new_03}. Since $K\geqslant 6\geqslant 4$ and $|S|\leqslant\eta n$, it suffices to show that
\[ 12 e^2 (d-1)^{L_0} (\eta/e)^{\frac{1}{2}}\leqslant\frac{1}{2}, \]
or, equivalently, that $\eta\leqslant 24^{-2}e^{-3}(d-1)^{-2L_0}$.
This follows from the choice of $\eta$ in \eqref{e-pap1-p1} and the fact that $d\geqslant 3$.

We proceed to the proof \eqref{e-pap-p5}. By \eqref{"douleia"}, it is enough to show that for every $\ell\in[T]$,
\begin{equation} \label{eq:new_05}
0.004 K \delta_\ell (d-1)^\ell |S|\geqslant \ell+1
\end{equation}
and
\begin{equation} \label{eq:new_06}
6e^2(d-1)^{L_0}\frac{(d-1)^\ell}{\delta_\ell} \Big(\frac{|S|}{en}\Big)^{1-\frac{6}{0.008 K \delta_\ell (d-1)^\ell}}
\leqslant\frac{1}{2}.
\end{equation}
To this end, set $m_\ell:= \log_{r}(\delta_\ell)$ for every $\ell\in[T]$.
\begin{claim}\label{claim_new_1}
 For every $\ell\in[T]$, we have $\ell - \frac{m_\ell}{2}\geqslant \frac{\ell}{2}-\frac{1}{2}$.
\end{claim}
\begin{proof}[Proof of Claim \ref{eq:new_06}]
Let $\ell\in[T]$. If $\ell\leqslant\lfloor 2T/3\rfloor$, then $m_\ell = \ell$. On the other hand,
if $\ell \geqslant \lfloor 2T/3\rfloor+1$, then $T\leqslant 3\ell/2$ and $m_\ell = 2T -2\ell +1$;
therefore, $\ell - \frac{m_\ell}{2} = 2\ell-T-\frac{1}{2}\geqslant \frac{\ell}{2}-\frac{1}{2}$, as desired.
\end{proof}
By Claim \ref{claim_new_1} and  \eqref{eq:new_04}, we have for every $\ell\in [T]$,
\begin{equation} \label{"douleia2"}
2^{\ell - \frac{m_\ell}{2}}\geqslant \frac{6}{0.008 K}(\ell+1)
\geqslant \max\Big\{ \frac{2}{0.008 K}(\ell+1), \frac{12}{0.008 K}\Big\}.
\end{equation}
\begin{claim}\label{claim_new_3}
For every $\ell\in[T]$, we have $\ell + \frac{m_\ell}{2} - T +\frac{6T}{0.008 K 2^{\ell-\frac{m_\ell}{2}}}\leqslant1$.
\end{claim}
\begin{proof}[Proof of Claim \ref{claim_new_3}]
Fix $\ell\in[T]$. By Claim \ref{claim_new_1}, it is enough to show that $\ell + \frac{m_\ell}{2} - T +\frac{6\sqrt{2}T}{0.008 K 2^{\frac{\ell}{2}}}\leqslant1$.

First, assume that $\ell \leqslant \lfloor T/2 \rfloor$. By \eqref{eq:new_04}, we see that $\frac{0.008K}{6}\geqslant 4$. Hence, using the fact that in this case $m_\ell=\ell$, we obtain that
\[ \ell + \frac{m_\ell}{2} - T +\frac{6\sqrt{2}T}{0.008 K 2^{\ell/2}}\leqslant-\frac{T}{4} + \frac{T}{0.008 K/6 }\leqslant0.\]

Next, assume that $\ell\in\{\lfloor T/2\rfloor+1,\dots,\lfloor2T/3\rfloor\}$. Then, $m_\ell = \ell$.
Since that function $f(x)=x2^{-\frac{x}{4}}$ is upper bounded by $2\sqrt{2}$, by \eqref{eq:new_04}, we have
\[ 1\geqslant \frac{6\sqrt{2}T}{0.008 K 2^{T/4}} \geqslant \frac{6\sqrt{2}T}{0.008 K 2^{\ell/2}}
\geqslant \ell +\frac{m_\ell}{2} -T + \frac{6\sqrt{2}T}{0.008 K 2^{\ell/2}}. \]

Finally, assume that $\ell \in \{\lfloor2T/3\rfloor+1,\dots,T\}$.
Then, $m_\ell = 2T - 2\ell +1$. Using, in this case, the fact that the function $f(x)=x2^{-\frac{x}{3}}$ is upper bounded by $\frac{7}{6}\sqrt{2}$, by \eqref{eq:new_04}, we conclude that
\[ 1\geqslant \frac{1}{2} + \frac{6\sqrt{2}T}{0.008 K 2^{T/3}} \geqslant
\frac{1}{2} + \frac{6\sqrt{2}T}{0.008 K 2^{\ell/2}}=
\ell +\frac{m_\ell}{2} -T + \frac{6\sqrt{2}T}{0.008 K 2^{\ell/2}}. \qedhere \]
\end{proof}
We are now ready to finish the proof.
First note that \eqref{eq:new_05} follows from \eqref{"douleia2"}
after observing that $\delta_\ell (d-1)^\ell \geqslant 2^{\ell-\frac{m_\ell}{2}}$.
In order to show \eqref{eq:new_06}, we fix $\ell\in [T]$. Since $d\geqslant3$, $(d-1)^T|S|\leqslant \eta n$
and $\delta_\ell^{-1} = 2^{\frac{m_\ell}{2}} \leqslant (d-1)^{\frac{m_\ell}{2}}$, by \eqref{"douleia2"} and
Claim \ref{claim_new_3}, we have
\begin{align*}
 6e^2&(d-1)^{L_0}\frac{(d-1)^\ell}{\delta_\ell} \Big(\frac{|S|}{en}\Big)^{1-\frac{6}{0.008 K \delta_\ell (d-1)^\ell}} \\
  & \leqslant 6e^2(d-1)^{L_0}(d-1)^{\ell+\frac{m_\ell}{2}-T+\frac{6T}{0.008 K \delta_\ell (d-1)^\ell}}
  \Big(\frac{\eta}{e}\Big)^{1-\frac{6}{0.008 K \delta_\ell (d-1)^{\ell}}} \\
  & \leqslant 6e^2(d-1)^{L_0}(d-1)^{\ell+\frac{m_\ell}{2}-T+\frac{6T}{0.008 K 2^{\ell-\frac{m_\ell}{2}}}}
  \Big(\frac{\eta}{e}\Big)^{1-\frac{6}{0.008 K (d-1)^{\ell-\frac{m_\ell}{2}}}} \\
  & \leqslant 6e^2(d-1)^{L_0+1} \Big(\frac{\eta}{e}\Big)^{1/2} \leqslant\frac{1}{2},
\end{align*}
where the last inequality holds by the choice of $\eta$ in \eqref{e-pap1-p1}. The proof of Proposition \ref{p3-pap1}
is thus completed.
\end{proof}

\end{document}
